\subsection{DEFINITION OF A UNIFORMITY BY MEANS OF A FAMILY OF PSEUDOMETRICS}

We have seen in Chapter VI, § 2, no. 3 that if, for each real number \(a > 0\) , we denote by \(\mathbf{U}_a\) the set of all pairs \((x, y)\) of points of \(\mathbf{R}^n\) whose Euclidean distance apart is \(\leqslant a\) , then the \(\mathbf{U}_a\) form a fundamental system of entourages of the uniformity of \(\mathbf{R}^n\) as \(a\) runs through the set of real numbers \(> 0\) .

More generally, let \(f\) be a pseudometric on a set \(\mathbf{X}\) ; for each \(a > 0\) , let \(\mathbf{U}_a\) denote \(\overline{f}^{1}([o, a])\) , and let us show that, as \(a\) runs through the set of all real numbers \(> 0\) , the \(\mathbf{U}_a\) form a fundamental system of entourages of a uniformity on \(\mathbf{X}\) . Axiom \((\mathbf{U}_{\mathbf{I}}^{\prime})\) is satisfied by reason of \((\mathrm{EC}_{\mathbf{I}})\) ; if \(a \leqslant b\) , we have \(\mathbf{U}_a \subset \mathbf{U}_b\) and therefore the \(\mathbf{U}_a\) satisfy \((\mathbf{B}_{\mathbf{I}})\) ; by \((\mathrm{EC}_{\mathbf{II}})\) , we have \(\overline{\mathbf{U}}_a = \mathbf{U}_a\) and therefore \((\mathbf{U}_{\mathbf{II}}^{\prime})\) is satisfied; finally, by \((\mathrm{EC}_{\mathbf{III}})\) we have \(\overline{\mathbf{U}}_a \subset \mathbf{U}_{2a}\) , so that \((\mathbf{U}_{\mathbf{III}}^{\prime})\) is satisfied. {[}Cf. Chapter II, § 1, no. 1, Definition 2{]}. Consequently we may make the following definition:

DEFINITION 2. Given a pseudometric \(f\) on a set \(\mathbf{X}\) , the uniformity defined by \(f\) is the uniformity on \(\mathbf{X}\) which has as a fundamental system of entourages the family of sets \(\overline{f}^{1}([0, a])\) , where a runs through the set of all real numbers \(> 0\) .

Two pseudometrics on X are said to be equivalent if they define the same uniformity.

Remarks. 1) If \((a_{n})\) is any sequence of numbers \(>0\) and tending to \(0\) , the \(\mathrm{U}_{a_n}\) form a fundamental system of entourages of the uniformity defined by \(f\) .

\begin{enumerate}
\def\labelenumi{\arabic{enumi})}
\setcounter{enumi}{1}
\tightlist
\item
  The definition of a uniformity by a pseudometric \(f\) consists in taking as a fundamental system of entourages the inverse image under \(f\) of the neighbourhood filter of \(o\) in the subspace \([o, +\infty]\) of \(\overline{\mathbf{R}}\) . Note that this procedure is quite analogous to that which allowed us to define the uniformities on a topological group (Chapter III, § 3, no. 1).
\end{enumerate}

Let \(f\) and \(g\) be two pseudometrics on \(\mathbf{X}\) . From Definition 2 it follows that the uniformity defined by \(f\) is coarser than the uniformity defined by \(g\) if and only if, for each \(a > 0\) there exists \(b > 0\) such that the relation \(g(x, y) \leqslant b\) implies \(f(x, y) \leqslant a\) . A necessary and sufficient condition for \(f\) and \(g\) to be equivalent pseudometrics is that for each \(a > 0\) there exists \(b > 0\) such that \(g(x, y) \leqslant b\) implies \(f(x, y) \leqslant a\) , and \(f(x, y) \leqslant b\) implies \(g(x, y) \leqslant a\) .

In particular, if there exists a constant \(k\) such that \(f \leqslant kg\) , the uniformity defined by \(f\) is coarser than the uniformity defined by \(g\) .

Let \(\varphi\) be a mapping of the interval \([o, +\infty]\) into itself, satisfying the following conditions: 1) \(\varphi(o) = o\) , and \(\varphi\) is continuous at \(o\) ; 2) \(\varphi\) is increasing in \([o, +\infty]\) and is strictly increasing in a neighbourhood of \(o\) ; 3) for all \(u \geqslant o\) and \(v \geqslant o\) , we have \(\varphi(u + v) \leqslant \varphi(u) + \varphi(v)\) . Then if \(f\) is any pseudometric on a set X, the composition \(g = \varphi \circ f\) is a pseudometric equivalent to \(f\) .

The reader may easily verify that we may, for example, take \(\varphi\) to be any one of the following functions:

\[
\sqrt {u}, \quad \log (\mathrm{I} + u), \quad \frac {u}{\mathrm{I} + u}, \quad \inf (u, \mathrm{I}).
\]

The last two examples show that there always exist bounded pseudometrics equivalent to any given pseudometric (finite or not).

DEFINITION 3. If \((f_{t})_{t\in I}\) is a family of pseudometrics on a set \(X\) , then the least upper bound of the set of uniformities defined on \(X\) by the pseudometrics \(f_{t}\) is called the uniformity defined by the family \((f_{t})\) .

Two families of pseudometrics on X are said to be equivalent if they define the same uniformity on X.

From the definition of the least upper bound of a set of uniformities (Chapter II, § 2, no. 5), the filter of entourages of the uniformity U defined on X by a family of pseudometrics \((f_{t})_{t\in\mathbf{I}}\) is the filter generated (Chapter I, §6, no. 2) by the family of sets \(\overline{f}_{t}^{1}([o,a])\) , where t runs through I and a runs through the set of real numbers \textgreater0. In other words, we obtain a fundamental system of entourages of U by proceeding as follows: we take at random a finite number of indices \(\iota_{1},\iota_{2},\ldots,\iota_{n}\) and, corresponding to each \(\iota_{k}\) , a number \(a_{k}>0\) ; then we consider the set V of pairs \((x,y)\in\mathbf{X}\times\mathbf{X}\) such that \(f_{\iota_{k}}(x,y)\leqslant a_{k}\) for \(i\leqslant k\leqslant n\) ; these sets V (for all possible choices of n, the \(\iota_{k}\) and the \(a_{k}\) ) form a fundamental system of entourages for U. Moreover, we may restrict ourselves to the case in which all the \(a_{k}\) are equal to the same number a\textgreater0, since the entourage consisting of all pairs \((x,y)\) such that

\[
\sup _ {1 \leqslant k \leqslant n} \left(f _ {\iota_ {k}} (x, y)\right) \leqslant \inf _ {1 \leqslant k \leqslant n} a _ {k}
\]

is evidently contained in V.

For each finite subset H of I, let \(g_{H}\) denote the upper envelope of the family \((f_{t})_{t\in\mathbb{H}}\) . As H runs through the set of all finite subsets of I and a runs through the set of real numbers \textgreater0, the sets \(g_{\mathrm{H}}^{-1}([o,a])\) form a fundamental system of entourages of the uniformity U. Now the \(g_{H}\) are pseudometrics on X (no. 1), and the upper envelope of a finite number of functions of the family \((g_{\mathrm{H}})\) belongs to this family, by definition; we express this property by saying that the family of pseudometrics \((g_{\mathrm{H}})\) is saturated. The family of pseudometrics \((g_{\mathrm{H}})\) is therefore equivalent to the family \((f_{t})\) , and is said to be the family of pseudometrics obtained by saturating \((f_{t})\) . From what has just been said it follows that we may always restrict ourselves to considering uniformities defined by saturated families of pseudometrics.

In the particular case where \(\mathbf{I}\) is a finite set, this argument shows that the uniformity defined by the family of pseudometrics \((f_{t})_{t\in \mathbf{I}}\) is also defined by the single pseudometric \(g = \sup_{t\in \mathbf{I}}f_t\) .

Let \(\mathcal{U},\mathcal{U}'\) be two uniformities on \(\mathbf{X}\) , defined respectively by two saturated families \((f_{t})_{t\in \mathbf{I}},(g_{x})_{x\in \mathbf{K}}\) . Then \(\mathcal{U}\) is coarser than \(\mathcal{U}'\) if and only if, for each index \(t\in I\) and each real number \(a > 0\) , there is an index \(x\in K\) and a number \(b > 0\) such that the relation \(g_{x}(x,y)\leqslant b\) implies \(f_{t}(x,y)\leqslant a\) .

Example of a uniformity defined by a family of pseudometrics. Let \((f_{t})_{t\in I}\) be an arbitrary family of (finite) real-valued functions defined on a set \(X\) . Let \(\mathcal{U}\) be the coarsest uniformity on \(X\) with respect to which the \(f_{t}\) are uniformly continuous (Chapter II, § 2, no. 3). Then it follows from the definition of the entourages of \(\mathcal{U}\) (loc. cit.) that \(\mathcal{U}\) is the uniformity defined on \(X\) by the pseudometrics

\[
g _ {i} (x, y) = | f _ {i} (x) - f _ {i} (y) |.
\]

